\begin{exercisebox}{RSL representation}
    Proof
    \begin{equation}
        C(z) = \left(\frac{\ln(1-z)}{1-z}\right)_+ \,\leftrightarrow\, C^R(z) = 0, C^S(z) = \frac{\ln(1-z)}{1-z}, C^L(x) = \frac{\ln(1-x)^2}{2}
    \end{equation}
\end{exercisebox}
\noindent{} Hint: read \cref{sec:rsl}.

\begin{exercisebox}{Product of functions 1}
    Give the RSL representation of the product a regular function $r(z)$ and a +-distribution $\qty(p(z))_+$, $C(z) = r(z)\qty(p(z))_+$
\end{exercisebox}
\noindent{} A popular example found in many text books is
\begin{equation}
    P_{qq}^{(0)}(y) = C_F\left( \frac{1+y^2}{(1-y)_+} + \frac 3 2 \delta(1-y)\right) \,\forall q \label{eq:Pqq}
\end{equation}
where the first term has the requested structure, i.e.\ the +-distribution is only intended for the denominator.

\begin{exercisebox}{Product of functions 2}
    Proof
    \begin{equation}
        \left[r(z)p(z)\right]_+ = r(z) \left[p(z)\right]_+ - \delta(1-z) \int\limits_0^1 dy~ r(y) \left[p(y)\right]_+
    \end{equation}
    where again $r(z)$ is a regular function and $\qty(p(z))_+$ is a +-distribution.
\end{exercisebox}
\noindent{} Check that \cref{eq:Pqq} can be also written as
\begin{equation}
    P_{qq}^{(0)}(y) = C_F\left( \frac{1+y^2}{1-y}\right)_+ \,\forall q\,.
\end{equation}
Another (lengthy) example is \cref{eq:NLOCq} vs.\ \cref{eq:NLOCqp}.
